\documentclass{beamer}

\usetheme{Berlin}
\usecolortheme{Orchid}
\useoutertheme{miniframes}
\setbeamertemplate{navigation symbols}{}

\usepackage[utf8]{inputenc}
\usepackage[T1]{fontenc}
\usepackage{amsmath}
\usepackage{amssymb}
\usepackage{booktabs}
\usepackage{graphicx}
\graphicspath{{../images/}}

\usepackage{xcolor}
\usepackage{listings}
\usepackage{hyperref}
\usepackage{tikz}
\usetikzlibrary{arrows.meta,positioning,shapes.geometric,calc}

\lstset{
  language=Java,
  basicstyle=\ttfamily\small,
  keywordstyle=\color{blue},
  commentstyle=\color{gray},
  stringstyle=\color{teal},
  showstringspaces=false,
  columns=fullflexible,
  breaklines=true
}

% Per-slide source line (references shown on the slide itself).
\newcommand{\slidesources}[1]{%
  \vfill
  {\tiny\color{gray}\raggedright\sloppy\parbox{\linewidth}{Sources: #1}\par}%
}

% Unit-wise source strings for this subject (used by slides/notes).
% Path is relative to the lecture `latex/` folder (the compilation CWD).
% OOPS with Java (CSEG1044) — unit-wise sources
%
% These are short, student-facing reference strings intended to be used with:
%   - \slidesources{...}  (in beamer slides)
%   - \notesources{...}   (in lecture notes)
%
% Style inspiration: other subjects in My-Subjects/ use semicolon-separated sources
% and include an "accessed" date for web documentation.

\newcommand{\OOPSUnitISources}{Schildt, \textit{Java: A Beginner's Guide} (10e), McGraw-Hill (Java fundamentals + OOP basics).; Oracle Java Tutorials: Getting Started + Language Basics + Classes and Objects (accessed \today).; Java SE API docs (e.g., \texttt{String}, \texttt{Scanner}) (accessed \today).}

\newcommand{\OOPSUnitIISources}{Schildt, \textit{Java: The Complete Reference} (12e), McGraw-Hill (classes, inheritance, packages, interfaces).; Bloch, \textit{Effective Java} (3e), Addison-Wesley, 2018 (design choices; interfaces vs abstract classes).; Oracle Java Tutorials: Classes and Objects + Interfaces + Packages (accessed \today).; Java SE API docs (e.g., \texttt{Comparable}, \texttt{Runnable}) (accessed \today).}

\newcommand{\OOPSUnitIIISources}{Schildt, \textit{Java: The Complete Reference} (12e) (nested/inner classes, exceptions, threads, I/O).; Oracle Java Tutorials: Nested Classes + Exceptions + Concurrency + Basic I/O (accessed \today).; Java SE API docs (e.g., \texttt{Thread}, \texttt{AutoCloseable}) (accessed \today).}

\newcommand{\OOPSUnitIVSources}{Schildt, \textit{Java: The Complete Reference} (12e) (generics, lambdas, Swing, JDBC).; Bloch, \textit{Effective Java} (3e) (generics + lambdas).; Oracle Java Tutorials: Generics + Lambda Expressions + Swing + JDBC Basics (accessed \today).; Java SE API docs (e.g., \texttt{java.util.function}, \texttt{javax.swing}, \texttt{java.sql}) (accessed \today).}

\newcommand{\OOPSUnitVSources}{Schildt, \textit{Java: The Complete Reference} (12e) (Collections Framework + wrapper classes).; Oracle Java docs: Collections Framework overview (accessed \today).; Java SE API docs (\texttt{java.util} collections; wrapper classes in \texttt{java.lang}) (accessed \today).}

\newcommand{\OOPSUnitVISources}{Git SCM: \textit{Pro Git} (2e) (version control workflow), accessed \today.; GitHub Docs (repositories, issues, pull requests), accessed \today.; Oracle Java Tutorials: Swing + JDBC (accessed \today).; JUnit 5 User Guide (accessed \today).; Apache Maven Project: Getting Started (accessed \today).; Gradle User Manual: Getting Started (accessed \today).; UML class diagrams: UML 2.x overview/spec (accessed \today).}

\newcommand{\OOPSMiscWhyInterfacesSources}{Schildt, \textit{Java: The Complete Reference} (12e) (interfaces).; Bloch, \textit{Effective Java} (3e) (prefer interfaces where appropriate).; Oracle Java Tutorials: Interfaces (accessed \today).; Java SE API docs (e.g., \texttt{Comparable}, \texttt{Runnable}) (accessed \today).}



\title[OOPS with Java]{Object Oriented Programming (CSEG1044)}
\subtitle{Unit 03 -- Lecture 01: Nested Classes (Types + Use-cases)}
\author{Tofik Ali}
\institute{School of Computer Science, UPES Dehradun}
\date{\today}

\begin{document}

\begin{frame}[fragile]
  \titlepage
  \vspace{0.5em}
  \slidesources{\OOPSUnitIIISources}
\end{frame}

\begin{frame}{Quick Links}
  \centering
  \hyperlink{sec:overview}{\beamerbutton{Overview}}\hspace{0.6em}
  \hyperlink{sec:concepts}{\beamerbutton{Key Topics}}\hspace{0.6em}
  \hyperlink{sec:demo}{\beamerbutton{Demo}}\hspace{0.6em}
  \hyperlink{sec:summary}{\beamerbutton{Summary}}
  \slidesources{\OOPSUnitIIISources}
\end{frame}

\begin{frame}{Agenda}
  \tableofcontents
  \slidesources{\OOPSUnitIIISources}
\end{frame}

\section{Overview}
\label{sec:overview}

\begin{frame}{Learning Outcomes}
  \begin{itemize}[<+->]
    \item Explain why nested classes improve encapsulation in some designs.
    \item Differentiate static nested, inner, local, and anonymous classes.
    \item Recognize use cases: helpers and callbacks/listeners (Swing later).
  \end{itemize}
  \slidesources{\OOPSUnitIIISources}
\end{frame}

\section{Key Topics}
\label{sec:concepts}

\begin{frame}{Unit 03: Nested Classes, Exceptions, Multithreading \& IO Streams}
  \begin{itemize}[<+->]
    \item Lecture focus: Nested Classes (Types + Use-cases)
  \end{itemize}
  \slidesources{\OOPSUnitIIISources}
\end{frame}

\begin{frame}{Today: Roadmap}
  \begin{itemize}[<+->]
    \item Nested class motivation: scope and encapsulation
    \item Static nested, inner, local, and anonymous classes
    \item Use-cases: helpers, callbacks/listeners (bridge to Swing later)
  \end{itemize}
  \slidesources{\OOPSUnitIIISources}
\end{frame}

\begin{frame}{Why nested classes?}
  \begin{itemize}[<+->]
    \item Keep helper types close to the code that uses them.
    \item Reduce public API surface (better encapsulation).
    \item Common for callbacks/listeners.
  \end{itemize}
  \slidesources{\OOPSUnitIIISources}
\end{frame}

\begin{frame}{Types (quick)}
  \begin{itemize}[<+->]
    \item \textbf{Static nested}: no outer instance reference.
    \item \textbf{Inner}: has implicit outer instance reference.
    \item \textbf{Local}: declared inside a method; limited scope.
    \item \textbf{Anonymous}: unnamed; used inline for one-time behavior.
  \end{itemize}
  \slidesources{\OOPSUnitIIISources}
\end{frame}

\begin{frame}{Static nested vs inner (key difference)}
  \begin{itemize}[<+->]
    \item \textbf{Static nested}: behaves like a normal class scoped inside outer; no outer object needed.
    \item \textbf{Inner}: every instance is tied to an outer object (has an implicit outer reference).
  \end{itemize}
  \vspace{0.6em}
  \textbf{Rule:} if you do not need outer state, prefer \texttt{static}.
  \slidesources{\OOPSUnitIIISources}
\end{frame}

\begin{frame}[fragile]{Example: static nested helper (validator)}
\begin{lstlisting}
class BankAccount {
  static class Validator {
    static void requirePositive(double amt) {
      if (amt <= 0) throw new IllegalArgumentException("amt");
    }
  }

  void deposit(double amt) {
    Validator.requirePositive(amt);
    // ... update balance
  }
}
\end{lstlisting}
  \slidesources{\OOPSUnitIIISources}
\end{frame}

\begin{frame}[fragile]{Example: local class (inside a method)}
\begin{lstlisting}
public void demoLocal() {
  final int times = 2; // effectively final

  class Printer {
    void say(String msg) {
      for (int i = 0; i < times; i++)
        System.out.println(msg);
    }
  }

  new Printer().say("Hi");
}
\end{lstlisting}
  \slidesources{\OOPSUnitIIISources}
\end{frame}

\begin{frame}{Use cases}
  \begin{itemize}[<+->]
    \item Helpers: validators, builders, entries/tuples
    \item Iterators and small adapters
    \item Callbacks/listeners (Swing event handling later)
  \end{itemize}
  \slidesources{\OOPSUnitIIISources}
\end{frame}


\section{Demo / Example}
\label{sec:demo}

\begin{frame}[fragile]{Example: inner + anonymous (small) (1/2)}
\begin{lstlisting}
interface Greeter { void greet(String name); }

public class Outer {
    private final String prefix = "Hello";

    class InnerGreeter implements Greeter {
        public void greet(String name) {
            System.out.println(prefix + ", " + name);
        }
    }
\end{lstlisting}
  \slidesources{\OOPSUnitIIISources}
\end{frame}

\begin{frame}[fragile]{Example: inner + anonymous (small) (2/2)}
\begin{lstlisting}
// continued...
    public void demo() {
        Greeter a = new InnerGreeter();
        Greeter b = new Greeter() { // anonymous
            public void greet(String name) {
                System.out.println(prefix + " (anon), " + name);
            }
        };
        a.greet("UPES");
        b.greet("UPES");
    }
}
\end{lstlisting}
  \slidesources{\OOPSUnitIIISources}
\end{frame}

\begin{frame}{Mini Exercises (2--3 mins)}
  \begin{enumerate}[<+->]
    \item Which type has an implicit outer reference?
    \item Which type is commonly used for callbacks?
    \item What does ``effectively final'' mean for local classes?
  \end{enumerate}
  \slidesources{\OOPSUnitIIISources}
\end{frame}


\section{Summary}
\label{sec:summary}

\begin{frame}{Key Takeaways}
  \begin{itemize}[<+->]
    \item Nested classes improve encapsulation and keep helpers local.
    \item Know the 4 kinds: static nested, inner, local, anonymous.
    \item Anonymous classes are common for callback-style code (Swing).
  \end{itemize}
  \slidesources{\OOPSUnitIIISources}
\end{frame}

\begin{frame}{Checkpoint Questions}
  \begin{enumerate}[<+->]
    \item Difference: static nested vs inner class?
    \item Why are anonymous classes used in event handling?
    \item Where can a local class be used?
  \end{enumerate}
  \slidesources{\OOPSUnitIIISources}
\end{frame}

\begin{frame}{Next Lecture Preview}
  \textbf{Next:} Exception Handling (Core)
  \vspace{0.6em}
  \begin{itemize}[<+->]
    \item We will handle errors safely using \texttt{try/catch/finally} and custom exceptions.
  \end{itemize}
  \slidesources{\OOPSUnitIIISources}
\end{frame}

\end{document}
